% ============================================================
%  ch03c  §3.3 志村簇 mod p 纤维上的一般 Tate 圈
%  原文对应：ICM-arXiv.tex 1144--1245 行
%  编号：定理 3.3.1；注 3.3.2；例 3.3.3；公式 (3.3)(3.4)
% ============================================================

\subsection{志村簇 mod $p$ 纤维上的一般 Tate 圈}\label{S: Tate}

在 \cite{Xiao-Zhu-Cycle} 中，我们应用定理 \ref{introT: spectral action} 验证了许多志村簇的 mod $p$ 纤维上 Tate 猜想的“一般”情形。沿用 \S\ref{S: Congruence} 的记号，并（为简单起见）保留其中的假设。以 $(\overline{\mathscr S}_{K,\overline{k_v}})^\pf$ 记 $\overline{\mathscr S}_{K,\overline{k_v}}$ 的完美化（即把它看作 $\Aff_{\overline{k_v}}^{\pf}$ 上的完美预层）；对 $\overline{\mathscr S}_{K,\bar{k}}$ 的每一点附上一个带 $G$-结构的 $F$-同晶体（见 \cite{Kisin-Integral-Model, Xiao-Zhu-Cycle}），便可定义所谓的 Newton 映射
\[
\on{Nt}: (\overline{\mathscr S}_{K,\overline{k_v}})^\pf\to \kot(G_{\bQ_p})_{\overline{k_v}}.
\]
于是 $\kot(G_{\bQ_p})_{\overline{k_v}}$ 的 Newton 分层（见 \S\ref{S: Local Langlands}）诱导出 $\overline{\mathscr S}_{K,\overline{k_v}}$ 被局部闭子簇分层。已知 $\on{Nt}$ 的像含唯一基本元 $b$，$\overline{\mathscr S}_{K,\overline{k_v}}$ 中相应的子簇是闭的，称为基本 Newton 层，记作 $\overline{\scrS}_b$。

设 $m$ 是几何 Frobenius $\sigma_p$ 在 $\xch(\hat{T})$ 上作用的阶。令
\[
\Lambda^\Tate_p=\Big\{\la\in \xch(\hat T)\; \Big|\; \sum_{i=0}^{m-1} \sigma_p^i(\la)=0\Big\}\subset \xch(\hat T).
\]
对 $\hat G_{\bQ_\ell}$ 的表示 $V$ 与 $\la\in\xch(\hat T)$，以 $V(\la)$ 记 $V$ 的 $\la$-权子空间（关于 $\hat T$），并令
\[
V^\Tate=\bigoplus_{\la\in\Lambda^\Tate_p} V(\la).
\]
我们特别关心条件 $V^\Tate_\mu\neq 0$。如 \cite{Xiao-Zhu-Cycle} 的引言所说明，在志村簇上同调的猜想性 Galois 理论描述（类似于猜想 \ref{C: gen coh}）之下，对一个在 $p$ 处的 Satake 参数足够一般的 Hecke 模 $\pi_f$，该向量空间维数的某个倍数 $a(\pi_f)$ 应等于 $\overline{\scrS}_{K,\overline{k_v}}$ 中间维数紧支撑上同调的 $\pi_f$-分支中 Tate 类空间的维数。此外，这个空间通常很大。例如，当 $G$ 是二次虚域上符号为 $(i,n-i)$ 的奇（射影）酉群时，在惰性素数处该空间的维数为 $\begin{pmatrix}\frac{n+1}{2}\\ i\end{pmatrix}$。

对代数闭域上维数为 $d$ 的（未必不可约）代数簇 $Z$，以 $H^{\on{BM}}_{2d}(Z)(-d)$ 记最高次 Borel--Moore 同调的 $(-d)$-Tate 扭，它是以 $Z$ 的不可约分支为基的向量空间。今设 $X$ 是 $q$ 元有限域 $k$ 上维数为 $d+r$ 的光滑簇，$Z \subseteq X_{\overline k}$ 是维数为 $d$ 的（未必不可约）射影子簇。有圈类映射
\[
\cl: H^{\on{BM}}_{2d}(Z)(-d)\to \bigcup_{j\geq 1}H_{c}^{2d}(X_{\overline k}, \bQ_\ell(d))^{\sigma_q^j}=:T^d_\ell(X).
\]

\begin{thm}
\label{T:main theorem}
记 $d_{\mu}=\dim X=2d$，$r=\dim V_\mu^{\Tate}$。
\begin{enumerate}
\item\label{T:main theorem-1} $\overline{\scrS}_{K,\overline{k_v}}$ 的基本 Newton 层 $\overline{\mathscr S}_b$ 是纯维数 $d$ 的，特别地 $d$ 总是整数。此外有 $H_{K,\bQ_\ell}$-等变同构
\[
H^{\on{BM}}_{2d}(\overline{\mathscr S}_{b})(-d)\cong C(G'(\bQ)\backslash G'(\bA_f)/K,\bQ_\ell)^{\oplus r},
\]
其中 $G'$ 是 $G$ 的有限平凡化内形式，且 $G'_\bR$ 紧。
\item\label{T:main theorem-2} 设 $\pi_f$ 是 $H_{K,\Ql}$ 的不可约模，令
\[
H^{\on{BM}}_{2d}(\overline{\mathscr S}_{b})[\pi_f]=\Hom_{H_{K,\Ql}}(\pi_f,H^{\on{BM}}_{2d}(\overline{\mathscr S}_{b})(-d)_{\Ql})\otimes \pi_f
\]
为 $\pi_f$-等型分支。则圈类映射
\[
\cl:H^{\on{BM}}_{2d}(\overline{\mathscr S}_{b})(-d)\to T^d_\ell(\overline{\mathscr S}_{K})
\]
限制在 $H^{\on{BM}}_{2d}(\overline{\mathscr S}_{b})[\pi_f]$ 上单射，只要 $\pi_{f,p}$（即 $\pi_f$ 在 $p$ 处的分支）的 Satake 参数是 $V_\mu$-一般的。
\item\label{T:main theorem-3} 设 $\Sh_K(G)$ 是（本质上的）四元数志村簇或 Kottwitz 算术簇。若 $\pi_{f,p}$ 的 Satake 参数是强 $V_\mu$-一般的，则圈类映射的 $\pi_f$-等型分支满射到 $T^d_\ell(\overline{\scrS_K})[\pi_f]$。特别地，对这些 $\pi_f$，Tate 猜想成立。
\end{enumerate}
\end{thm}

\begin{rmk}
\begin{enumerate}
\item 对 $\hat G$ 的表示 $V$，“$V$-一般”与“强 $V$-一般” Satake 参数的定义见 \cite[Definition 1.4.2]{Xiao-Zhu-Cycle}。$\CG|_{d=(p,\sigma_p)}$ 中的正则半单元总是 $V$-一般的，但反之不然。见 \cite[Remark 1.4.3]{Xiao-Zhu-Cycle}。
\item 该定理的一些特殊情形最初在 \cite{Helm-Tian-Xiao,Tian-Xiao-Tate} 中证明。
\end{enumerate}
\end{rmk}

该定理的证明依赖几个不同的要素。借助志村簇基本部分的 Rapoport--Zink 单值化，第 \eqref{T:main theorem-1} 部分可以化归为关于某些仿射 Deligne--Lusztig 簇的不可约分支的问题，后者在 \cite[\S 3]{Xiao-Zhu-Cycle} 中研究过。最难的是第 \eqref{T:main theorem-2} 部分，我们的证明是计算 $\overline\scrS_b$ 中所有 $d$ 维圈之间的相交数。这些数可以编码为一个元素在 $H^{\cl}_{K_p}$ 中的 $r\times r$ 矩阵。一般说来，直接具体地算出这个矩阵似乎没有指望。

然而，这个矩阵可以理解为 $\Coh(\Loc_{\CG,p}^{\ur})$ 中某些态射的复合。具体地说，先把 $G'(\bQ)\backslash G'(\bA)/K$ 实现为志村集，其志村余特征为 $\mu'=0$。然后利用定理 \ref{introT: spectral action}（以及 Satake 同构 \eqref{E: derived Satake}），该矩阵可计算为
\[
\begin{aligned}
&\Hom_{\Coh(\Loc_{\CG,p}^{\ur})}(\mO,\widetilde{V_{\mu}})\otimes  \Hom_{\Coh(\Loc_{\CG,p}^{\ur})}(\widetilde{V_{\mu}},\mO)\\
&\qquad\to  \Hom_{\Coh(\Loc_{\CG,p}^{\ur})}(\mO,\mO)\cong H^{\cl}_{K_p}\otimes\bQ_\ell.
\end{aligned}
\]
接着需要判定这一配对何时非退化，这本身是表示论中一个有意思的问题，其解决依赖对向量值函数的 Chevalley 限制映射的研究。该矩阵的行列式在 \cite{Xiao-Zhu-Invariants} 中算出。最后，第 \eqref{T:main theorem-3} 部分通过比较两个迹公式（关于 $G$ 的 Lefschetz 迹公式与关于 $G'$ 的 Arthur--Selberg 迹公式）得到证明。

\begin{ex}\label{Tate cycle on U(1,2n)}
设 $G=\on{U}(1,2r)$ 是二次虚扩张 $E/\bQ$ 所关联的 $(2r+1)$ 元酉群\footnote{它不是伴随群，故本例与我们的假设不尽相符；但这样讨论起来更方便，计算本质上也一样。}，其在无穷处的符号为 $(1,2r)$。它带有标准的志村数据，从而（在固定水平 $K\subset G(\bA_f)$ 之后）给出志村簇。特别地，若 $r=1$，这就是（本质上的）Picard 模曲面。设 $p$ 是在 $E$ 中惰性的素数，且 $K_p$ 超特殊。此时 $\overline{\scrS}_b$ 是某些 Deligne--Lusztig 簇之并，由 $G'(\bQ)\backslash G'(\bA_f)/K$ 参数化，其中 $G'=\on{U}(0,2r+1)$ 在所有有限位置同构于 $G$。这些圈在 $\overline{\scrS}_b$ 内部的相交模式（本质上）已由 \cite{Vollaard-Wedhorn-GU(n1)} 给出，但圈之间的相交数要难算得多。事实上，除把定理 \ref{introT: spectral action} 应用于此情形之外，我们并不知道对一般的 $r$ 如何直接计算它们。（$r=1$ 的情形可以直截了当地处理。）

此时 $\hat G=\GL_{2r+1}$，$\sigma_p$ 在其上按 $A\mapsto J(A^T)^{-1}J$ 作用，其中 $J$ 是反对角线上元素全为 $1$ 的反对角矩阵。表示 $V_\mu$ 是 $\GL_{2r+1}$ 的标准表示。可以验证 $\dim V_\mu^{\Tate}=1$（这与上面所说 $\overline{\scrS}_b$ 的不可约分支由 $G'(\bQ)\backslash G'(\bA_f)/K$ 参数化一致）。按惯例把 $\hat G$ 的权格等同于 $\bZ^{2r+1}$。此时 $\Hom_{\Coh(\Loc_{\CG,p}^{\ur})}(\mO,\widetilde{V_{\mu}})$ 是 $\Hom_{\Coh(\Loc_{\CG,p}^{\ur})}(\mO,\mO)=H_{K_p}^{\cl}\otimes\bQ_\ell$ 上的自由秩一 module。于是生成元 $\mathbf a_{\on{in}}$ 诱导出 $H_{K,\bQ_\ell}$-等变同态
\[S(\mathbf a_{\on{in}}): C(G'(\bQ)\backslash G'(\bA_f)/K)\to H_c^{2r}(\overline{\scrS}_{K,\overline{k_v}},\bQ_\ell(r)),\]
它实现 $\overline\scrS_b$ 的圈类映射（差一个倍数）。模 $\Hom_{\Coh(\Loc_{\CG,p}^{\ur})}(\widetilde{V_{\mu}},\mO)$ 同样是 $H_{K_p,\bQ_\ell}$ 上的自由秩一模。对所选生成元 $\mathbf a_{\on{out}}$，复合
\[S(\mathbf a_{\on{out}})\circ S(\mathbf a_{\on{in}})=S(\mathbf a_{\on{out}}\circ\mathbf a_{\on{in}})\]
从 $\overline\scrS_b$ 的不可约分支计算出这些圈的相交矩阵。

元素 $h:=\mathbf a_{\on{out}}\circ\mathbf a_{\on{in}} \in H_{K_p,\bQ_\ell}$ 已在 \cite[Example 6.4.2]{Xiao-Zhu-Invariants} 中显式计算（差显然的修正，并借助 Satake 同构 \eqref{Sat isom classical}）。具体地，
\begin{equation}
\label{E:Hecke operator in Tpj}
h=p^{r(r+1)}\sum_{i=0}^{r} (-1)^i(2i+1)p^{(i-r)(r+i+1)}\sum_{j=0}^{r-i} \begin{bmatrix} 2r+1-2j \\ r-i- j\end{bmatrix}_{t=-p}T_{p,j}.
\end{equation}
这里 $T_{p,j}=1_{K_p\la_j(p)K_p}$，其中 $\la_i=(1^i,0^{2r-2i+1},(-1)^i)$；而 $\begin{bmatrix} n \\ m\end{bmatrix}_t$ 是二项式系数的 $t$-类比，由
\[
[0]_t=1,\quad [n]_t=\frac{t^n-1}{t-1},\quad [n]_t!=[n]_t[n-1]_t\cdots[1]_t,\quad \begin{bmatrix} n \\ m\end{bmatrix}_t=\frac{[n]_t!}{[n-m]_t![m]_t!}
\]
给出。换言之，此时 $\overline\scrS_b$ 中圈的相交矩阵由 Hecke 算子 \eqref{E:Hecke operator in Tpj} 算出。

这一计算有一个有意思的推论，即对有限 Deligne--Lusztig 簇的相交理论的一个结论，而我们并不知道它的直接证明。设 $W$ 是 $\bF_{p^2}$ 上 $(2r+1)$ 维非退化 Hermite 空间。考虑 $r$ 维 Deligne--Lusztig 簇
\[
\mathrm{DL}_{r}: = \big\{ H \subset W \textrm{ 维数为 } r\;\big|\; H \subseteq (H^{(p)})^\perp \big\},
\]
其中 $H^{(p)}$ 是 $H$ 沿 Frobenius 的拉回。以 $\mH$ 记相应的秩 $r$ 万有子丛。令 $\mE = \mH^{(p)} \otimes \big( (\mH^{(p)})^\perp / \mH \big)$。则
\begin{equation}\label{E: chern class on DL}
\int_{\mathrm{DL}_r} c_r( \mE)  = \sum_{i = 0}^r (-1)^i(2i+1) p^{i^2+i} \begin{bmatrix} 2r+1 \\ r-i\end{bmatrix}_{t=-p}.
\end{equation}
\end{ex}
